% The abstract of the notes.  Three parts, in this order: Översikt,
% Lärandemål, Förkunskaper.
%
% The objectives reuse the wording of the lecture decks' objectives where
% an exercise practises exactly that objective (Upprepningar's IterLO*,
% Behållare: Listor's ListLO*, Behållare: Tupler's TupleLO* and Moduler
% och paket's PkgLO*), because \cref cannot cross documents: an \ltnote
% here has to point at a label in this document.  Only the first one,
% \cref{OvningIterationsLOReview}, is the tutorial's own -- no lecture
% deck claims it, but the course's outcome Lo11 does.

\emph{Översikt}
Den här övningen använder föreläsningarnas innehåll i stället för att visa
det.  Vi börjar med ett program som någon annan skrivit och som inte gör
vad det ser ut att göra, och letar upp felen i det.  Därefter byggs en
frågesport upp, en aspekt i taget: frågorna läggs i en lista och gås
igenom, antalet försök avgör om det blir \mintinline{python}{for} eller
\mintinline{python}{while}, en meny låter användaren hålla på så länge hen
vill, det som flera program behöver flyttar ut i en egen modul, och till
sist blandas frågornas ordning med en funktion som bara står i
dokumentationen.  Sedan följer fyra fördjupande uppgifter: den sista står
för sig själv, medan de tre första har en given ordning sinsemellan.
Varje uppgift följs av ett lösningsförslag --- ett av flera möjliga ---
så att du kan fortsätta på egen hand efteråt.

\emph{Lärandemål}
Efter övningen ska studenten kunna:
\begin{restatable}{lo}{OvningIterationsLOReview}%
  \label{OvningIterationsLOReview}%
  Läsa ett program som någon annan har skrivit, hitta vad som är fel
  eller kan bli bättre, och rätta det. % Lo11, Lo02
\end{restatable}
\begin{restatable}{lo}{OvningIterationsLOFor}\label{OvningIterationsLOFor}%
  Använda \mintinline{python}{for} för att göra samma sak en gång för
  varje element i en samling. % Lo04
\end{restatable}
\begin{restatable}{lo}{OvningIterationsLOWhile}%
  \label{OvningIterationsLOWhile}%
  Använda \mintinline{python}{while} för att upprepa så länge ett villkor
  gäller, och se till att villkoret kan bli falskt. % Lo04
\end{restatable}
\begin{restatable}{lo}{OvningIterationsLORecursion}%
  \label{OvningIterationsLORecursion}%
  Läsa och skriva rekursiva funktioner med basfall och rekursivt anrop på
  ett mindre problem. % Lo04
\end{restatable}
\begin{restatable}{lo}{OvningIterationsLOChoose}%
  \label{OvningIterationsLOChoose}%
  Välja mellan \mintinline{python}{for}, \mintinline{python}{while} och
  rekursion utifrån hur problemet ser ut, och motivera valet. % Lo04, Lo02
\end{restatable}
\begin{restatable}{lo}{OvningIterationsLOInput}%
  \label{OvningIterationsLOInput}%
  Skriva en återanvändbar funktion för inmatning som frågar om igen tills
  användaren har matat in data som går att använda. % Lo08, Lo01, Lo03
\end{restatable}
\begin{restatable}{lo}{OvningIterationsLOList}%
  \label{OvningIterationsLOList}%
  Gå igenom en lista, element för element och med både elementets nummer
  och elementet. % Lo09, Lo04
\end{restatable}
\begin{restatable}{lo}{OvningIterationsLOAppend}%
  \label{OvningIterationsLOAppend}%
  Lägga till element i en lista medan programmet kör, och avgöra om en
  ändring som görs inne i en funktion syns hos den som anropade
  den. % Lo09
\end{restatable}
\begin{restatable}{lo}{OvningIterationsLOTuple}%
  \label{OvningIterationsLOTuple}%
  Packa upp en tupel i flera variabler, och använda det för att låta en
  funktion lämna tillbaka flera värden. % Lo09, Lo03
\end{restatable}
\begin{restatable}{lo}{OvningIterationsLOModule}%
  \label{OvningIterationsLOModule}%
  Lägga det som flera program behöver i en egen modul och hämta det
  därifrån med \mintinline{python}{import}. % Lo03, Lo01
\end{restatable}
\begin{restatable}{lo}{OvningIterationsLOMain}%
  \label{OvningIterationsLOMain}%
  Skriva en modul så att dess provkörning bara körs när filen körs
  direkt, och inte när ett annat program hämtar den. % Lo03, Lo04
\end{restatable}
\begin{restatable}{lo}{OvningIterationsLODocs}%
  \label{OvningIterationsLODocs}%
  Hitta och läsa dokumentationen för en modul eller en behållare, och
  därifrån använda något som inte gåtts igenom på
  föreläsningen. % Lo06, Lo09
\end{restatable}

\emph{Förkunskaper}
Studenten bör ha tagit del av föreläsningarna \emph{Upprepningar},
\emph{Behållare: Listor}, \emph{Behållare: Tupler} och \emph{Moduler och
paket}: \mintinline{python}{for}, \mintinline{python}{while} och
rekursion, listor och tupler med index och \mintinline{python}{len}, samt
\mintinline{python}{import} och villkoret som skiljer en importerad modul
från en fil som körs direkt.  Därtill förutsätts egna funktioner med
parametrar och returvärde från föreläsningen \emph{Funktioner}, och
\mintinline{python}{try}/\mintinline{python}{except} kring
\mintinline{python}{int} från föreläsningen \emph{Inmatning och
felhantering}.
